% =====================================================================
\part{Os alicerces}
% =====================================================================

\chapter{O dado: o que é, exatamente, um pixel do MapBiomas}

A espinha dorsal do trabalho é a Coleção 10.1 do MapBiomas, série anual de uso
e cobertura da terra com resolução de 30 metros, de 1985 a 2024. Entender o que
essa frase significa é a base de tudo o mais, porque três propriedades do dado
governam o que se pode e o que não se pode medir com ele.

\section{Anual, e não instantâneo}

Cada ano da série é uma classificação, e não uma fotografia. O produto atribui
a cada célula de 30 por 30 metros uma classe (formação florestal, formação
savânica, pastagem, agricultura, área urbana e assim por diante) a partir de um
conjunto de imagens do ano inteiro, processado por algoritmo. Duas
consequências decorrem daí.

A primeira é que a série permite construir a \textbf{matriz de transição}: para
cada par de anos, a contagem de quantos pixels saíram da classe $i$ e chegaram
à classe $j$. É isso que autoriza dizer que um hectare de lavoura \emph{veio}
de pastagem, e não apenas que a lavoura cresceu enquanto a pastagem encolheu.
Nenhuma diferença entre estoques permite essa afirmação.

A segunda é que o classificador tem erro, e o erro dele não é aleatório no
tempo nem no espaço. Duas manifestações desse erro condicionam o trabalho
inteiro, e são tratadas explicitamente no texto.

\section{A oscilação entre pastagem e savana}

Nas matrizes de transição aparece um fluxo reverso volumoso, de pastagem para
vegetação natural, que seria a evidência natural de regeneração. Examinado, ele
se revela outra coisa: de 75\% a 98\% dele, conforme o par de anos, dirige-se à
formação \emph{savânica} --- ou seja, à classe mais parecida com pasto sujo aos
olhos do classificador. E a razão entre o fluxo de ida e o de volta colapsa de
8,4 vezes, no primeiro par decenal, para 1,22 vez no par anual mais recente.

O diagnóstico é o seguinte: regeneração real é lenta e unidirecional; oscilação
de classificador é bidirecional e aproximadamente balanceada. O que se observa
tem a segunda assinatura.

\begin{naodiz}
Isso \emph{não} significa que não haja regeneração alguma em Goiás. Há um
componente real, minoritário, restrito à formação florestal, que só se acumula
em horizontes longos. O que a constatação proíbe é narrar o fluxo bruto de
pastagem para vegetação natural como recuperação.
\end{naodiz}

\begin{nabanca}
``Esse fluxo reverso é grande e eu não o leio como regeneração. Ele é quase
todo dirigido à formação savânica e é aproximadamente simétrico ao fluxo de
ida, o que é assinatura de oscilação de classificador na borda entre pasto e
cerrado, e não de recuperação. As conclusões centrais descansam no fluxo direto
de conversão e em saldos líquidos, que são robustos a uma oscilação simétrica
porque ela se cancela no líquido.''
\end{nabanca}

\section{A migração de rótulo do Mosaico de Usos}

Esta é a peça mais importante do trabalho inteiro, e vale entendê-la a fundo.

No trecho final da série, parte da conversão de pastagem em lavoura passa a
receber o rótulo da classe 21, ``Mosaico de Usos'', em vez de ``Agricultura''.
O efeito na leitura é dramático: pela classe ``Agricultura'' isolada, a
conversão de pasto em lavoura no Sul goiano \emph{cai 88\%} entre o Ato~II e o
Ato~III. Pela mesma janela, a área de soja que o IBGE levanta em campo
\textbf{cresce 38\%} no estado. As duas medidas não podem estar certas ao mesmo
tempo sobre o mesmo território.

A resolução é a decisão D26, e ela é uma decisão de honestidade e não de
conveniência. Toda medida de agricultura em janela que alcance 2020--2024 é
reportada como um \textbf{intervalo} entre duas réguas:

\begin{itemize}
\item o \textbf{piso}, que é a classe ``Agricultura'' isolada, e que
\emph{subconta}, porque perde a lavoura reetiquetada;
\item o \textbf{teto}, que é a união de ``Agricultura'' com ``Mosaico de
Usos'', e que \emph{superconta}, porque o Mosaico contém também sistemas
integrados de lavoura e pecuária reais, mosaicos finos demais para separar a
30 metros, e mosaico antigo e estável, anterior ao fenômeno.
\end{itemize}

Considera-se robusta apenas a conclusão que sobrevive nos dois extremos.

\begin{armadilha}[A tentação de tratar a união como ``a medida certa'']
A união não é correção: somá-las suporia que todo o Mosaico é agricultura mal
rotulada, o que é falso. Ela responde a uma pergunta \emph{diferente e mais
grossa}: ``quanta terra saiu de pastagem para lavoura ou uso misto?''. Se a
banca sugerir que você adote a união e siga em frente, essa é a resposta.
\end{armadilha}

Três evidências mostram que o erro de medida aponta \emph{contra} a tese, e não
a favor dela --- ponto que vale ter na ponta da língua, porque é o que converte
uma fragilidade aparente em argumento.

\begin{enumerate}
\item Entre 2019 e 2024, quatro medidas indicam deslocamento ao norte (três
delas entre 10 e 13 km, a quarta de 4,4 km) e a quinta indica que nada se moveu
--- justamente a única cujo rótulo mudou no período. Rebanho e soja plantada,
medidos em campo, marcham.
\item Localizada no mapa, a área que mudou de rótulo não está espalhada: o
centro dessa massa fica \textbf{46 km ao norte} do centro da agricultura
visível. A conversão que o satélite deixou de nomear está acontecendo na
fronteira.
\item O crescimento por região dessa massa acompanha o da soja do IBGE quase
ponto a ponto: correlação de 0,84, e 1,52 contra 1,53 milhão de hectares.
\end{enumerate}

\begin{nabanca}
``A régua crua faz a marcha parecer \emph{mais fraca} do que foi. Se eu quisesse
inflar o resultado, bastaria adotar a união e não declarar nada. O que fiz foi o
contrário: reportei o intervalo e ancorei numa fonte que não passa pelo
classificador.''
\end{nabanca}

\section{Onde a União vale e onde ela não vale}

Uma distinção fina, e que a banca pode testar. A régua da união serve para uma
janela curta (2019--2024), onde a migração de rótulo de fato ocorre. Aplicada
aos quarenta anos, ela devolve $-60$ km, isto é, a agricultura andando para o
\emph{sul}, e esse número não significa o que parece.

A razão é que em 1985 a classe Mosaico já existia e era extensa (3,63 Mha,
10,7\% do estado), situada bem ao norte, e não era lavoura sob outro nome: era
terra de uso misto indiferenciado. Ao longo da série ela encolhe enquanto a
lavoura cresce ao sul, de modo que somar as duas classes nas duas pontas é
comparar objetos diferentes. A confirmação é espacial: o Mosaico ocupa
praticamente a mesma área nas duas pontas da série, mas as manchas quase não se
sobrepõem, e a correlação espacial entre as células é de apenas \textbf{0,15}.
O de 1985 se dissolveu; o de 2024 surgiu em outro lugar.

\chapter{As unidades: 246 municípios e 166 AMCs}

Esta é a decisão de rigor mais consequente do trabalho, e a que um economista
na banca vai reconhecer imediatamente.

\section{O problema}

A malha municipal brasileira não é constante no tempo, e Goiás é um caso
severo: o estado tinha 184 municípios com estatística municipal em 1985 e tem
246 hoje. Dos 246, \textbf{62 (25\%)} só aparecem na estatística depois de
1985, em ondas de emancipação concentradas em 1989 (27), 1993 (21), 1997 (10) e
2001 (4).

O que torna isso fatal é que as duas fontes se comportam de maneiras
\emph{opostas} diante de uma emancipação:

\begin{itemize}
\item o \textbf{raster} é recortado pelo polígono \emph{atual} em todos os
anos, de modo que a série de cobertura do município-pai já vem no território
reduzido de hoje, inclusive para 1985;
\item a \textbf{estatística do IBGE} é tabulada pelo município \emph{como ele
existia naquele ano}, de modo que a produção do pai em 1985 inclui o território
que depois virou o filho.
\end{itemize}

Ao juntar as duas pelo código municipal, compara-se a cobertura do território de
hoje contra a produção de um território historicamente maior. E o efeito não é
sutil: municípios-pai individuais registram quedas de 50\% a 80\% no rebanho
\emph{no ano da emancipação}, enquanto o total estadual atravessa esses anos
sem descontinuidade. Ler essas quedas como dinâmica pecuária seria atribuir a
um fenômeno econômico o que é perda de território, e como as ondas de 1993 e
1997 caem dentro das janelas das análises longitudinais, elas dominariam
qualquer estimador de variação no período.

\section{A solução e o seu preço}

A agregação em \textbf{Áreas Mínimas Comparáveis} reúne cada município-pai aos
seus desmembramentos numa unidade de território constante ao longo de toda a
janela. Aplicada a Goiás com a concordância de 1980--2010, janela que antecede
todas as emancipações do período analisado, ela reduz os 246 municípios a
\textbf{166 AMCs}: 53 grupos de pai e filhos e 113 unidades que permanecem
idênticas ao município.

A validação é direta: o total estadual é idêntico ao da malha municipal, e
nenhuma das 166 AMCs estreia depois de 1985.

\begin{ancora}[A aritmética da AMC, dita com precisão]
246 municípios $\to$ 166 AMCs. Os 53 grupos reúnem \textbf{133 municípios}, que
deixam de ser analisáveis isoladamente nas séries longas. A malha perde
\textbf{80 unidades} em relação às 246.
\end{ancora}

\begin{naodiz}
A AMC não recupera informação: ela impede uma comparação inválida. A informação
histórica do município-filho \emph{não existe} separada, porque estava no pai. O
preço é inevitável quando o dado vem de pesquisa.
\end{naodiz}

\section{A regra de agregação, que é onde os erros nascem}

Ao agregar, grandezas \textbf{extensivas} (hectares, cabeças, reais) são
somadas. Grandezas \textbf{derivadas} (razões, densidades, lotação) são
\emph{recalculadas} a partir das extensivas já somadas, e nunca somadas nem
promediadas diretamente, o que produziria média de razões em lugar de razão
de totais. Formalmente, para uma lotação $L$:

\begin{equation}
L_{\text{AMC}} = \frac{\sum_{m \in \text{AMC}} \text{cabeças}_m}{\sum_{m \in \text{AMC}} \text{hectares}_m}
\qquad\text{e nunca}\qquad
L_{\text{AMC}} \neq \frac{1}{n}\sum_{m \in \text{AMC}} \frac{\text{cabeças}_m}{\text{hectares}_m}.
\end{equation}

\section{A estratégia de dois trilhos}

Usar o trilho errado é erro metodológico, e o trabalho declara qual é qual:

\begin{itemize}
\item \textbf{166 AMCs} para toda análise \emph{longitudinal}: primeiras
diferenças, painel com efeitos fixos, periodização, tendências, deslocamento de
centros de massa.
\item \textbf{246 municípios} para análises \emph{transversais} e para o
período recente: o Censo Agropecuário de 2017, os mapas de um único ano, a
conta de desenvolvimento com o IFDM.
\end{itemize}

\section{A lotação em cabeças por hectare, e não em Unidade Animal}

Detalhe pequeno com resposta pronta, porque um zootecnista na banca vai
estranhar. A conversão para Unidade Animal aplica um fator constante ao efetivo
do rebanho. Por ser \emph{constante}, ele se cancela em qualquer análise de
tendência ou de variação, que é onde a lotação entra neste trabalho. Por ser
\emph{fixo no tempo}, ele produziria um nível não comparável entre 1985 e 2024,
já que o peso médio do animal mudou ao longo do período. A medida reportada é,
portanto, a diretamente observada.

\chapter{O painel e a primeira diferença}

\section{A tabela}

As fontes convergem para uma tabela única, organizada por município e ano:
$246 \times 40 = 9.840$ linhas e 195 colunas. É desse painel, e de sua versão
agregada em AMCs, que sai toda a estimação do trabalho.

\section{Por que nada é estimado em nível (decisão D7)}

Esta é a decisão estruturante, e há duas razões independentes para ela.

\textbf{Razão 1: tendência comum.} Quase toda série do painel cresce ao longo de
quatro décadas --- área, produção, produto, população. Correlacionar
\emph{níveis} entre séries com tendência comum produz associações altas e
vazias: elas medem o fato de que ambas cresceram, e nada mais.

\textbf{Razão 2: identidade contábil disfarçada.} Muitas variáveis do painel são
internamente relacionadas por construção. Área plantada e produção, ou área e
valor adicionado, medem facetas do mesmo fato e apresentam correlação em nível
da ordem de \textbf{0,9}. Estimar uma sobre a outra em nível mede uma
identidade, e não uma relação.

A solução é diferenciar:
\begin{equation}
\Delta x_t = x_t - x_{t-1}.
\end{equation}

E a reformulação da pergunta é o ponto que vale dizer em voz alta. Em vez de
\emph{``municípios com mais crédito têm menos pasto?''}, que é uma comparação
entre lugares diferentes e confunde tudo o que difere entre eles, o modelo
passa a responder \emph{``quando o crédito de um município subiu, o pasto dele
recuou?''}.

\begin{ancora}[A auditoria que sustenta a decisão]
As mesmas duplas que têm correlação em nível da ordem de 0,9 apresentam
correlação próxima de \textbf{zero} em variação \emph{within}. A primeira
diferença, combinada com os efeitos fixos, desarma o vazamento.
\end{ancora}

\section{As três exceções declaradas}

A regra governa o painel de associação. Três leituras do texto correm fora
dela, cada uma por exigência do próprio método, e saber nomeá-las evita a
acusação de inconsistência.

\begin{enumerate}
\item A \textbf{correlação em nível estadual} do quadro das camadas é
descritiva, e vem com os seus limites declarados na mesma linha em que aparece.
\item O procedimento de \textbf{Toda-Yamamoto} ajusta o vetor autorregressivo em
níveis \emph{por construção}, e é justamente isso que torna a sua inferência
válida sob integração (Capítulo~\ref{cap:granger}).
\item As equações do \textbf{teto de oferta} retiram a tendência por efeitos
fixos de unidade e de ano em vez de por diferença, porque o que elas medem é o
que acontece a uma unidade \emph{enquanto ela se esgota}, e diferenciar
apagaria justamente esse nível.
\end{enumerate}

\section{Por que 38, e não 40}

Convenção de contagem que aparece em toda parte e que a banca pode cobrar. A
série de cobertura tem 40 observações anuais (1985--2024). A primeira diferença
consome a primeira observação; a defasagem de um ano consome a segunda. Daí as
\textbf{38 diferenças anuais} que aparecem nas regressões, no cálculo de erros
robustos a autocorrelação e na contagem de realizações independentes do choque
macroeconômico.

\begin{nabanca}
``Quando eu digo 40, o objeto é a série descritiva. Quando digo 38, é o painel
estimado.''
\end{nabanca}

\chapter{Duas convenções geodésicas, e por que são duas}

Ponto técnico curto, com resposta que impressiona pela precisão.

Toda operação \textbf{geométrica} --- centroides, distâncias, deslocamentos,
recortes zonais --- é feita em \textbf{EPSG:5880}, que é o SIRGAS 2000 em
projeção Policônica do Brasil, padrão cartográfico brasileiro e a projeção dos
mapas do texto.

Mas a policônica é projeção de \emph{compromisso}: não é conforme nem
equivalente. Por isso ela \textbf{não é usada para medir área}. As áreas vêm de
outro caminho: o censo de pixels conta células de 30 m na grade nativa do
produto e converte a contagem em hectares pelo cosseno da latitude média
observada de cada município, sem passar por projeção nenhuma.

\begin{nabanca}
``Uso duas convenções porque elas respondem a perguntas diferentes. A policônica
me dá geometria comparável e é o padrão cartográfico do país; ela não preserva
área, e por isso área nenhuma é medida nela. A contagem de área é feita na grade
nativa do raster, com correção pelo cosseno da latitude.''
\end{nabanca}
